%% governance-deontic-logic: Verified Dyadic Deontic Logic and Governance Reasoning in Lean 4
%%
%% Render with:
%%   pdflatex governance-deontic-logic.tex && bibtex governance-deontic-logic \
%%   && pdflatex governance-deontic-logic.tex && pdflatex governance-deontic-logic.tex

\documentclass[]{article}
\usepackage[utf8]{inputenc}
\usepackage{url}
\usepackage[hyperindex,breaklinks]{hyperref}
%\usepackage{breakurl}  % disabled: incompatible with pdflatex/tectonic
\usepackage{listings}
\lstset{
  basicstyle=\ttfamily\footnotesize,
  breaklines=true,
  frame=single,
  extendedchars=false,
  literate={->}{$\to$}{2} {<->}{$\leftrightarrow$}{3}
           {forall}{$\forall$}{6} {exists}{$\exists$}{6}
           {<|}{{$\langle$}}{1} {|>}{{$\rangle$}}{1}
}
\usepackage{float}
\restylefloat{table}
\usepackage{longtable}
\usepackage{graphicx}
\usepackage[font=small,labelfont=bf]{caption}
\usepackage{amsfonts}
\usepackage{amssymb}
\usepackage{amsmath}
\usepackage{amsthm}

\newtheorem{theorem}{Theorem}
\newtheorem{lemma}[theorem]{Lemma}
\newtheorem{definition}[theorem]{Definition}
\newtheorem{remark}[theorem]{Remark}
\newtheorem{corollary}[theorem]{Corollary}

\newcommand{\code}[1]{\texttt{#1}}
\newcommand{\OB}{\ensuremath{\mathbf{O}}}
\newcommand{\PE}{\ensuremath{\mathbf{P}}}
\newcommand{\OP}{\ensuremath{\mathbf{Opt}}}
\newcommand{\DDL}{\textup{DDL+}}
\newcommand{\DTS}{\textup{DTS}}
\newcommand{\PLN}{\textup{PLN}}
\newcommand{\PGC}{\textup{PGC}}

\title{Verified Dyadic Deontic Logic and Governance Reasoning in Lean~4\\
  \texttt{(DRAFT --- \today)}}
\author{Zar\footnote{Corresponding author.}
\and Oru\v{z}i\footnote{``Oru\v{z}i'' denotes AI-assisted collaboration
using Claude Code (Anthropic, Opus 4.6), Codex (OpenAI, GPT-5.*), and
ChatGPT Pro (GPT-5.*). Co-authored under human direction.}}

\begin{document}
\maketitle

\begin{abstract}
We present a Lean~4 formalization of governance deontic reasoning,
spanning from classical deontic logic through eventuality reification to
constructive governance models and ethical foundations.
The formalization includes:
(i)~a shallow semantic embedding of Carmo and Jones' Dyadic Deontic Logic (\DDL{})
with five semantic conditions;
(ii)~a proof of all 13 Carmo-Jones axiomatic characterisation theorems (CJ\_3--CJ\_15)
as consequences of the embedding;
(iii)~the Deontic Traditional Scheme (\DTS{}) with obligation as the sole primitive,
derived permission and optionality, and all 12 von Wright interdefinability relations
as theorems;
(iv)~a governance reasoning layer with Hobbs-style eventuality reification, ISO~24617-4
thematic roles, a deontic query encoder, and a three-level judgment pipeline;
(v)~a constructive 2-state reactive loop (\code{GovNormCycle}) providing a provably
seriality-closed accessibility context;
(vi)~a temporal-deontic bridge formalizing obligation persistence under time translation
and a causal-normative prohibition theorem ruling out spurious causal candidates from
obligation;
(vii)~an abstract \code{ModalQueryEncoder} unifying event-calculus and deontic encoders
via a shared interface;
(viii)~an application to Gewirth's Principle of Generic
Consistency (\PGC{}); and
(ix)~formal refinement theorems connecting MeTTa governance rule implementations
to the Lean \DTS{} semantics, via a shared \code{MeTTaLike} evaluator abstraction
covering both PeTTa and HE~MeTTa evaluation, with proven \code{let*} unfolding
equivalences and kernel-checked canonical HE interpreter conformance tests.
We further formalize subsumption-based compliance between abstract and concrete
eventualities, and illustrate its use with a minimal proof-carrying treaty
kernel for AI-to-AI subcontracting.
Finally,~(x) we introduce an evidence-graded actuality layer that replaces the
metaphysical ``really exists'' reading of \code{rexist} with a two-layer
semantics: graded eventuality evidence and policy-relative categorical
acceptance.  A provenance-aware accepted-trace mechanism gates treaty
compliance on attestor trust and evidence quality, with a proven
soundness theorem linking the existing \code{RexistBridge} to
acceptance equivalence for any policy.
\end{abstract}

\tableofcontents
\newpage

%% ============================================================
\section{Introduction}
\label{sec:intro}
%% ============================================================

Formal governance and policy reasoning require a robust logical foundation that
can represent obligations, permissions, and temporal dynamics.
Classical propositional deontic logic
\cite{vonWright1951} provides a starting point, but real governance scenarios
involve:
\begin{itemize}
\item \emph{Contrary-to-duty obligations} --- norms that apply precisely when
      a primary obligation is violated \cite{CarmoJones2002};
\item \emph{Eventuality reification} --- treating events and states as
      first-class entities so that thematic roles and temporal structure
      can be reasoned about explicitly \cite{Hobbs1985};
\item \emph{Temporal dynamics} --- obligations that persist, shift, or
      expire as time advances; and
\item \emph{Causal spuriousness} --- distinguishing genuine causal obligations
      from extensional correlates that lack intensional support.
\end{itemize}

This paper describes a Lean~4 formalization in the \code{Mettapedia} repository
that addresses all four concerns.
The architecture layers from abstract to concrete:
Carmo-Jones \DDL{} semantics $\to$ \DTS{} algebra $\to$ governance reasoning layer
$\to$ reactive accessibility model $\to$ temporal and causal bridges $\to$
Gewirth \PGC{} application.

\paragraph{Contributions.}
The primary contributions are:
\begin{enumerate}
\item \textbf{Lean~4 port of the Benzm{\"u}ller et al.\ Isabelle/HOL DDL+ embedding}
      \cite{Benzmuller2018}, making it available in the Lean~4 ecosystem with
      Mathlib compatibility.  All 13 Carmo-Jones theorems are proven from
      5 semantic conditions (§\ref{sec:ddl}--\ref{sec:cj}).
\item \textbf{A Lean~4 port of Benzm{\"u}ller et al.'s Isabelle/HOL proof of
      Gewirth's Principle of Generic Consistency} (\PGC{})
      \cite{Benzmuller2019PGC}, connecting it to the DDL+ and \DTS{} layers
      above (§\ref{sec:pgc}).
\item \textbf{An honest architectural exploration} of whether formal deontic logic
      can ground governance reasoning, including a negative result:
      simple terminating languages do \emph{not} satisfy the seriality
      requirement needed for governance accessibility (§\ref{sec:govnorm}).
\item \textbf{A temporal-deontic bridge} formalizing obligation persistence
      under time translation and a causal-normative prohibition theorem
      connecting \PLN{} causal profiles to \DTS{} (§\ref{sec:temporal}).
\item \textbf{Formal MeTTa-to-Lean refinement theorems} proving that MeTTa
      governance rule implementations, when interpreted through a deontic
      pattern decoder, correspond to the proven \DTS{} theorems ---
      covering PeTTa and HE~MeTTa evaluators via a shared typeclass,
      with kernel-checked canonical HE interpreter conformance
      (§\ref{sec:refinement}).
\end{enumerate}

\noindent
Supporting contributions include: a \DTS{} algebra with derived permission
and optionality (§\ref{sec:dts}); a governance reasoning layer with Hobbs
eventuality reification, ISO~24617-4 thematic roles, and a three-level
verdict pipeline (§\ref{sec:govlayer}); an abstract \code{ModalQueryEncoder}
unifying event-calculus and deontic query dispatch (§\ref{sec:mqe}); and
a shared \code{MeTTaLike} evaluator typeclass with proven \code{let*}
unfolding equivalences (§\ref{sec:refinement}).

%% ============================================================
\section{Shallow Embedding of DDL+}
\label{sec:ddl}
%% ============================================================

\noindent
\textbf{Source attribution.}
The shallow-embedding architecture in this section --- including the
two-dimensional \code{WProp}/\code{Meaning} type structure, the
\code{DDLPlusFrame} with its five semantic conditions, the Kaplanian context
setup, and the LD-validity notion --- is a Lean~4 port and abstraction of the
Isabelle/HOL formalization by Benz{\"u}ller, Farjami, and Parent
\cite{Benzmuller2018}.  The Isabelle source files \code{CJDDLplus.thy} and
\code{ExtendedDDL.thy} are the direct predecessors of the Lean types defined
here.  Our contribution is the port to Lean~4 and the integration with the
\DTS{} governance layer.

\subsection{Types and Semantic Foundation}

The \DDL{} embedding follows the standard \emph{possible-worlds semantics}
tradition, representing propositions as truth-sets over a world type.
We adopt Kaplan's two-dimensional semantics as used by Benz{\"u}ller et al.\
\cite{Benzmuller2018,Kaplan1989} to handle context-dependent propositions
(``characters'').

\begin{definition}[World-Proposition and Meaning]
A \emph{world-proposition} over worlds \code{w} is a predicate
\code{WProp w}~$=$ \code{w -> Prop}.
A \emph{meaning} (Kaplan's character) is a context-dependent proposition
\code{Meaning c w}~$=$ \code{c -> w -> Prop}.
\end{definition}

These correspond to Isabelle types \texttt{wo = w => bool} and
\texttt{m = c => w => bool} in the source formalization \cite{Benzmuller2018}.

\subsection{The DDL+ Frame}

The semantic foundation of \DDL{} is a \emph{frame} bundling two accessibility
relations and a dyadic obligation relation:

\begin{lstlisting}
structure DDLPlusFrame (w : Type*) where
  av : w -> w -> Prop   -- actual-version accessibility
  pv : w -> w -> Prop   -- possible-version accessibility
  ob : WProp w -> WProp w -> Prop  -- dyadic obligation
  sem_3a : forall v, exists v', av v v'      -- av is serial
  sem_4a : forall v v', av v v' -> pv v v'   -- av <= pv
  sem_4b : forall v, pv v v                  -- pv is reflexive
  sem_5a : forall X, not (ob X (fun _ => False))
  sem_5b : forall X Y Z, (forall v, (X v /\ Y v) <-> (X v /\ Z v)) ->
             (ob X Y <-> ob X Z)
  sem_5c : forall X Y Z, (exists v, X v /\ Y v /\ Z v) ->
             ob X Y -> ob X Z -> ob X (fun v => Y v /\ Z v)
  sem_5d : forall X Y Z, (forall v, Y v -> X v) -> ob X Y ->
             (forall v, X v -> Z v) ->
             ob Z (fun v => (Z v /\ not (X v)) \/ Y v)
  sem_5e : forall X Y Z, (forall v, Y v -> X v) -> ob X Z ->
             (exists v, Y v /\ Z v) -> ob Y Z
\end{lstlisting}

\noindent
The five conditions \code{sem\_3a}--\code{sem\_5e} match Carmo \& Jones (2002)
conditions (3a), (4a), (4b), (5a)--(5e) \cite{CarmoJones2002}, pp.~290ff., and
correspond to lines 52--59 of the Isabelle source \code{CJDDLplus.thy}.

\begin{remark}
The accessibility relation \code{av} models ``open alternatives'' (what can
actually happen), while \code{pv} models ``possible alternatives'' (what could
ideally happen).  Condition \code{sem\_4a} ensures open alternatives are a
subset of possible ones, and \code{sem\_4b} ensures reflexivity of \code{pv}.
\end{remark}

\subsection{Modal and Deontic Operators}

Six modal operators are defined over \code{DDLPlusFrame}:

\begin{lstlisting}
def box_a (phi : Meaning c w) : Meaning c w :=  -- actual necessity
  fun ctx v => forall v', F.av v v' -> phi ctx v'
def box_p (phi : Meaning c w) : Meaning c w :=  -- possible necessity
  fun ctx v => forall v', F.pv v v' -> phi ctx v'
def box_S5 (phi : Meaning c w) : Meaning c w :=  -- unrestricted
  fun ctx _ => forall v, phi ctx v
\end{lstlisting}

\noindent
Diamond operators are defined dually.
Three deontic operators are then derived:

\begin{lstlisting}
-- Conditional obligation O<phi|sigma>
def cond_obl (phi sigma : Meaning c w) : Meaning c w :=
  fun ctx _ => F.ob (sigma ctx) (phi ctx)

-- Actual obligation: ob relative to av, with violation witness
def actual_obl (phi : Meaning c w) : Meaning c w :=
  fun ctx v => F.ob (F.av v) (phi ctx)
    /\ exists v', F.av v v' /\ not (phi ctx v')

-- Ideal obligation: ob relative to pv, with violation witness
def ideal_obl (phi : Meaning c w) : Meaning c w :=
  fun ctx v => F.ob (F.pv v) (phi ctx)
    /\ exists v', F.pv v v' /\ not (phi ctx v')
\end{lstlisting}

\noindent
The violation-witness condition (the $\exists v'$ clause) is standard in
contrary-to-duty frameworks: an obligation is ``ideal'' only when there exist
possible worlds where it is not satisfied.

\subsection{Kaplanian Context and LD Validity}

To model indexical statements, we follow the two-dimensional structure used by
Benz{\"u}ller et al.\ \cite{Benzmuller2018} (adapting Kaplan \cite{Kaplan1989}):

\begin{lstlisting}
structure KaplanianContext (c w e : Type*) where
  agent : c -> e    -- agent of the context
  world : c -> w    -- world of the context
\end{lstlisting}

\noindent
\emph{LD validity} (truth in every context at its own world) is strictly weaker
than classical modal validity.
This distinction is proven as a theorem:
\code{modal\_implies\_ld} shows that modal validity entails LD validity, but not conversely.
The \emph{a priori necessity} operator $\Box^D$ is truth at all LD contexts:
\code{box\_D phi ctx := ld\_valid phi}.

%% ============================================================
\section{Carmo-Jones Axiomatic Characterisation}
\label{sec:cj}
%% ============================================================

Carmo and Jones \cite{CarmoJones2002} present 13 axioms (CJ\_3--CJ\_15) as the
axiomatic characterisation of \DDL{}.  Following Benzm{\"u}ller et al.\
\cite{Benzmuller2018}, we derive all 13 from the 5 semantic frame conditions,
confirming the soundness of the axiomatic characterisation with respect to
the possible-worlds semantics.
Table~\ref{tab:cj} lists each theorem with its Lean identifier and the
primary frame condition used in its proof.

\begin{longtable}{lll}
\caption{Carmo-Jones theorems proven semantically in Lean~4}\label{tab:cj}\\
\textbf{Theorem} & \textbf{Lean identifier} & \textbf{Key condition} \\
\hline
CJ\_3: $\Box_p \Rightarrow \Box_a$ & \code{CJ\_3} & \code{sem\_4a} \\
CJ\_4: $\neg O\langle \bot | A \rangle$ & \code{CJ\_4} & \code{sem\_5a} \\
CJ\_5$^-$: conjunction of obligations & \code{CJ\_5\_minus} & \code{sem\_5c} \\
CJ\_6: obligation closed under conjunction & \code{CJ\_6} & \code{sem\_5e} \\
CJ\_7: modal equivalence preservation & \code{CJ\_7} & \code{sem\_5e} \\
CJ\_8: conditional equivalence preservation & \code{CJ\_8} & \code{C\_6} \\
CJ\_9a: $\Diamond_p O \Rightarrow \Box_p O$ & \code{CJ\_9a} & $O$ is world-independent \\
CJ\_9p: $\Diamond_a O \Rightarrow \Box_a O$ & \code{CJ\_9p} & $O$ is world-independent \\
CJ\_9\_var\_a: $O \Rightarrow \Box_a O$ & \code{CJ\_9\_var\_a} & definitional \\
CJ\_9\_var\_b: $O \Rightarrow \Box_p O$ & \code{CJ\_9\_var\_b} & definitional \\
CJ\_10: restriction under possibility witness & \code{CJ\_10} & \code{sem\_5e} \\
CJ\_11a\_var: conjunction of actual obligations & \code{CJ\_11a\_var} & \code{sem\_5c, sem\_3a} \\
CJ\_15: commitment extension & \code{CJ\_15} & \code{sem\_5d, sem\_5e} \\
\hline
\end{longtable}

\begin{remark}
CJ\_7 and CJ\_8 hold under \emph{classical modal validity} but not under the
weaker LD validity.  The Lean proofs make this explicit: they take a hypothesis
\code{modal\_valid (pequ A B)} rather than its LD-valid version, matching the
standard caveat in Carmo \& Jones (2002), §5.
\end{remark}

Additionally, the bridge between conditional, actual, and ideal obligation is
established in \code{DDLPlus/DTSBridge.lean}:
\code{axiomD\_actual} shows that actual obligation satisfies the KD axiom
(seriality of obligation), and \code{DDLPlusFrame\_to\_DTS} extracts a
\DTS{} structure from any \DDL{} frame.

%% ============================================================
\section{Deontic Traditional Scheme}
\label{sec:dts}
%% ============================================================

The Deontic Traditional Scheme (\DTS{}) \cite{AlchourrónBulygin1971,vonWright1951}
provides a three-valued classification of propositions: obligatory, permitted,
or optional.
We formalize it with a single primitive \code{ob} and derive the rest:

\begin{lstlisting}
structure DTS (P : Type*) where
  ob        : P -> Prop         -- obligation (primitive)
  neg       : P -> P            -- negation on propositions
  neg_neg   : forall p, neg (neg p) = p
  consistent : forall p, ob p -> not (ob (neg p))

-- Derived modalities
def DTS.pe (d : DTS P) (p : P) : Prop := not (d.ob (d.neg p))
def DTS.op (d : DTS P) (p : P) : Prop :=
  not (d.ob p) /\ not (d.ob (d.neg p))
\end{lstlisting}

\noindent
With these definitions, all 12 von Wright interdefinability relations become
theorems (Table~\ref{tab:dts}).

\begin{longtable}{ll}
\caption{DTS interdefinability theorems in Lean~4}\label{tab:dts}\\
\textbf{Statement} & \textbf{Lean identifier} \\
\hline
$\OB(p) \Rightarrow \PE(p)$ & \code{ob\_implies\_pe} \\
$\PE(p) \Leftrightarrow \neg\OB(\neg p)$ & \code{pe\_iff\_not\_ob\_neg} \\
$\OP(p) \Leftrightarrow \neg\OB(p) \wedge \neg\OB(\neg p)$ & \code{op\_iff} \\
$\neg\PE(p) \Rightarrow \OB(\neg p)$ & \code{not\_pe\_implies\_ob\_neg} \\
$\neg\OP(p) \Leftrightarrow \OB(p) \vee \OB(\neg p)$ & \code{not\_op\_iff} \\
$\PE(p) \wedge \PE(\neg p) \Leftrightarrow \OP(p)$ & \code{pe\_and\_pe\_neg\_iff\_op} \\
$\OB(p) \Rightarrow \neg\OP(p)$ & \code{ob\_implies\_not\_op} \\
$\OP(p) \Rightarrow \PE(p)$ & \code{pe\_of\_op} \\
$\OP(p) \Rightarrow \PE(\neg p)$ & \code{pe\_neg\_of\_op} \\
Trichotomy: $\OB(p) \vee \OP(p) \vee \OB(\neg p)$ & \code{dts\_trichotomy} \\
Exclusion: at most one of the three holds & \code{dts\_exclusive} \\
$\OB(p) \Rightarrow \neg\PE(\neg p)$ & \code{ob\_implies\_not\_pe\_neg} \\
\hline
\end{longtable}

\begin{remark}[On the nature of DTS theorems]
These interdefinability relations are largely definitional consequences of
\code{pe} and \code{op} being defined in terms of \code{ob} and \code{neg}.
The genuine content of the \DTS{} formalization is not any individual theorem
but rather the \emph{structure} itself: that obligation, permission, and
optionality form a consistent trichotomy under the consistency axiom
$\forall p,\, \OB(p) \to \neg \OB(\neg p)$.
The definitions do the mathematical work; the theorems confirm the
definitions are well-chosen.
\end{remark}

\noindent
The bridge \code{DDLPlusFrame\_to\_DTS} recovers a \DTS{} structure from a
\DDL{} frame by instantiating \code{ob} with the frame's dyadic obligation
restricted to the universal context and the reflexive world.

\begin{remark}[DTS consistency is not free]
\DTS{} consistency ($\OB(p) \to \neg\OB(\neg p)$) is \emph{not} derivable
from the \DDL{} frame conditions alone.  The \DDL{} axiom \code{sem\_5a}
prevents $\text{ob}(X, \bot)$ but does \emph{not} prevent simultaneous
$\text{ob}(X, p)$ and $\text{ob}(X, \neg p)$ when $X \cap p \cap \neg p = \emptyset$
(which it always is).  The \code{frameToDTS} bridge therefore requires an
explicit consistency hypothesis \code{hcons}, matching the Isabelle development
where \DTS{} is an independent axiomatisation.
\end{remark}

%% ============================================================
\section{Governance Reasoning Layer}
\label{sec:govlayer}
%% ============================================================

\subsection{Hobbs Eventuality Reification}

Following Hobbs \cite{Hobbs1985}, we treat events, states, and actions as
first-class entities:

\begin{lstlisting}
structure Eventuality (Entity Pred : Type*) where
  predicate : Pred                         -- e.g., soaMoor, soaPay
  roles     : ThematicRole -> Option Entity -- ISO 24617-4 assignments
  polarity  : Bool := true                 -- positive or negated
\end{lstlisting}

\noindent
The \code{ThematicRole} type enumerates 15 roles from ISO~24617-4:2014
\cite{ISO246174}: agent, patient, theme, experiencer, beneficiary, instrument,
location, time, source, goal, cause, result, manner, purpose, condition.
Polarity enables negated eventualities without duplicating the type:
\code{Eventuality.negate e} flips \code{e.polarity}, and
\code{negate\_negate} is a theorem (\code{!(!e) = e}).

\subsection{Deontic Query Encoder}

The \code{DeonticQueryEncoder} structure bridges the governance type system
and any world-model evidence calculus:

\begin{lstlisting}
structure DeonticQueryEncoder (Entity Pred Query : Type*) where
  modalQuery  : DeonticModality -> Eventuality Entity Pred -> Query
  groundQuery : CTTriple Entity Pred -> Query
\end{lstlisting}

\noindent
The five \code{DeonticModality} constructors are: \code{rexist} (alethic,
Hobbs' ``really exists''), \code{obligatory}, \code{permitted},
\code{forbidden} ($\OB(\neg p)$), and \code{optional} ($\neg\OB(p) \wedge \neg\OB(\neg p)$).
The \code{classify} function maps a proposition to its deontic status:
obligatory if $\OB(p)$, forbidden if $\OB(\neg p)$, optional otherwise.
The \code{modalQuery} field is the key interface: for a given modality and
eventuality it produces the WM query whose evidence represents the epistemic
degree to which that modality holds.

\subsection{Rexist Bridge and Evidence-Graded Actuality}
\label{sec:actuality}

The ``really exists'' (\code{rexist}) modality implements Hobbs' principle
that what really exists actually holds~\cite{Hobbs1985}.  In the Lean formalization,
\code{RexistBridge} is a \emph{hypothesis}: a \code{WMQueryEq} assumption
stating that rexist-evidence equals ground-evidence for the same predicate.
The bridge soundness theorem (\code{rexistBridgeRule}) simply unpacks
this assumption.  No modal axioms are needed --- the rexist bridge is an
interface contract between the governance layer and whatever world-model
evidence calculus is plugged in.

\paragraph{From ``really exists'' to eventuality evidence.}
The name \code{rexist} suggests binary metaphysical existence, but the
evidence it extracts is graded: a \PLN{} evidence pair $(n^+, n^-)$,
not a Boolean.  We therefore reinterpret \code{rexist} as an
\emph{eventuality-occurrence evidence channel}: ``how much evidence is
there that this eventuality obtained?''  The file
\code{ActualityPolicy.lean} introduces thin wrappers making this reading
explicit:
\begin{itemize}
\item \code{eventualityEvidence}~(alias: \code{eventEvidence}) ---
      extracts graded evidence that an eventuality occurred;
\item \code{projectionEvidence}~(alias for \code{groundEvidence}) ---
      extracts evidence for a structured factual projection (CT-triple);
\item \code{AcceptancePolicy} --- a predicate on evidence determining
      when graded evidence is sufficient for categorical governance
      acceptance.
\end{itemize}

\noindent
Categorical ``actuality for governance'' is then a downstream
\emph{acceptance policy}, not certainty-1 truth.  This follows the
formal-belief literature's separation of graded belief from categorical
acceptance, avoiding the lottery paradox that arises from naive
fixed-threshold approaches~\cite{Foley1993,Hawthorne2009}.

\paragraph{Three bridge strengths.}
The existing \code{RexistBridge} (evidence equality) is the strongest
possible bridge between the occurrence and projection channels.  We
introduce a weaker \code{AcceptanceBridge}: for a given policy~$\pi$,
occurrence is accepted if and only if the corresponding projection is
accepted, in every world-model state.  The key soundness theorem is:

\begin{theorem}[Acceptance bridge soundness]
\label{thm:acceptance-bridge}
For any acceptance policy~$\pi$, if the \code{RexistBridge} holds
(evidence equality), then the \code{AcceptanceBridge} holds.
\end{theorem}

\noindent
This shows the existing strong bridge is conservative over the weaker
acceptance-level bridge.  Three bridge strengths form a hierarchy:
evidence equality (strongest) $\Rightarrow$ strength agreement
(intermediate, future work) $\Rightarrow$ acceptance agreement (weakest,
most operationally realistic).

\paragraph{Support/confidence preorder.}
The raw coordinatewise order on evidence
($e_1 \le e_2 \Leftrightarrow n^+_1 \le n^+_2 \wedge n^-_1 \le n^-_2$)
is unsuitable for strength-based acceptance, since more evidence in that
order includes more negative evidence.  We define a parameterized
\emph{support/confidence preorder}:
\[
  e_1 \preceq_{(\alpha,\kappa)} e_2
  \;\;\Longleftrightarrow\;\;
  \mathrm{strength}_\alpha(e_1) \le \mathrm{strength}_\alpha(e_2)
  \;\wedge\;
  \mathrm{confidence}_\kappa(e_1) \le \mathrm{confidence}_\kappa(e_2)
\]
where $\alpha$ is a prior context and $\kappa$ a confidence parameter.
This is one member of a \emph{family} of preorders (reflexivity and
transitivity are proved), not a single canonical order on evidence.

\subsection{Three-Level Judgment Pipeline}

The \code{Judgments.lean} file defines a three-level inference pipeline:

\begin{enumerate}
\item \textbf{Eventuality level}: evidence that an eventuality exists with
      a given modality (\code{EventualityJudgment}).
\item \textbf{Statement level}: lifting eventuality judgments via derivation
      constructors --- Rexist bridge, DTS derivation, negation inference
      (\code{StatementJudgment}).
\item \textbf{Verdict level}: a \code{GovernanceVerdict} combining multiple
      statement judgments into an overall governance outcome.
\end{enumerate}

\noindent
The pipeline mirrors the three reasoning levels of the
\code{governance-reasoning-engine} source \cite{GovReasoningEngine}
(eventuality\_level, statement\_level, conclusion\_level).

\subsection{Subsumption-Based Compliance and Treaty Kernel}

The verdict layer detects contradiction, conflict, and violation using exact
event identity (\code{sameEvent}): same predicate and same role assignment.
This is appropriate for low-level consistency checking, but too rigid for
practical compliance, where an abstract norm is often satisfied by a more
specific concrete event.

We add a pure event-subsumption relation \code{eventSubsumes}, requiring
agreement on predicate and polarity together with role inclusion: every
thematic role specified in the abstract eventuality must be matched by the
concrete one.  This is the Lean counterpart of the MeTTa engine's
\code{is\_complied\_with\_by} relation, factored as a pure event relation.
To preserve backward compatibility, statement identities are introduced by a
wrapper \code{IdStatementJudgment} rather than by modifying the
\code{StatementJudgment} inductive.  The existing \code{governanceAnalysis}
remains exact-match; subsumption is introduced locally in a new treaty layer.

On top of event subsumption we define a minimal \emph{TreatyKernel} for
bilateral agent agreements.  A treaty clause specifies an obligor, obligee,
modality, eventuality pattern, optional deadline, and optional repair-clause
pointer.  The kernel distinguishes obligation fulfilment from prohibition
violation via separate predicates (\code{obligationFulfilled},
\code{prohibitionViolated}), avoiding the mistake of treating a forbidden
action as ``clause satisfaction.''  A concrete three-clause
confidential-subcall treaty (delivery deadline, nondisclosure,
refund-on-breach) demonstrates that more specific runtime events can discharge
abstract obligations and trigger breach detection.

\paragraph{Provenance-aware accepted traces.}
The treaty kernel operates over raw event traces, but real governance
requires distinguishing verified events from unreliable reports.
\code{TreatyKernelAcceptance.lean} introduces an
\code{AssessedTreatyEvent} wrapper pairing each treaty event with a
\PLN{} evidence assessment and attestor identity (using the existing
\code{attestedBy} field on \code{TreatyEvent}).  A provenance-aware
acceptance predicate filters the raw trace to an \emph{accepted base
trace}: only events from trusted attestors are admitted.

In the concrete subcall example, four candidate events arrive:
\begin{enumerate}
\item a self-claimed delivery by the specialist ($\langle 2,2\rangle$,
      attested by ``specialist'') --- \textbf{rejected};
\item a hash-verified delivery confirmed by an escrow oracle
      ($\langle 8,0\rangle$, ``hash-verifier'') --- \textbf{accepted};
\item a disclosure detected by a network monitor
      ($\langle 5,1\rangle$, ``monitor'') --- \textbf{accepted};
\item an unverified rumor from a gossip channel
      ($\langle 1,4\rangle$, ``gossip'') --- \textbf{rejected}.
\end{enumerate}

\noindent
Treaty compliance is then checked on the accepted trace:
the verified delivery fulfils the deliver-by-10 obligation, and the
monitor-detected leak violates the nondisclosure prohibition.
All acceptance decisions, the trace computation, and the treaty
compliance theorems are kernel-checked (\code{rfl}-proofs for
acceptance decisions, direct construction for compliance).  This
demonstrates the clean architectural separation: the evidence/attestation
layer decides which events are \emph{admitted}; the treaty layer
stays crisp and classical over the admitted trace.

\subsection{Worked Example: Rental Obligation}
\label{sec:rental}

We illustrate the full pipeline with a concrete scenario.

\paragraph{Setup.}
A rental agreement states:
\begin{enumerate}
\item[(P)] \emph{Primary obligation}: The tenant ought to pay rent.
  \hfill $\OB(\text{pay})$
\item[(CTD)] \emph{Contrary-to-duty}: If the tenant does not pay,
  the tenant ought to give notice.
  \hfill $\OB(\text{notice} \mid \neg\text{pay})$
\end{enumerate}

\noindent
This is a classic Chisholm-paradox scenario \cite{Chisholm1963}.
\DDL{}'s dyadic obligation handles it without contradiction because
$\OB(\text{pay})$ and $\OB(\text{notice} \mid \neg\text{pay})$ live at
different conditioning contexts.

\paragraph{Level 1: Eventuality judgment.}
The eventuality ``Alice pays Bob rent'' is encoded as:
\begin{lstlisting}
def e_pay : Eventuality Entity Pred :=
  { predicate := soaPay,
    roles := fun
      | .agent => some Alice
      | .patient => some Bob
      | .theme => some Rent
      | _ => none,
    polarity := true }
\end{lstlisting}
A world-model query \code{modalQuery .obligatory e\_pay} returns evidence
$\langle s, n \rangle$ (e.g., $s=5$ positive observations out of $n=7$ total),
yielding:
\begin{lstlisting}
EventualityJudgment .obligatory e_pay evidence
\end{lstlisting}

\paragraph{Level 2: Statement judgment.}
The eventuality judgment is lifted to a tagged statement referencing
the lease clause:
\begin{lstlisting}
StatementJudgment.fromL1 eventualityJudgment
\end{lstlisting}

\paragraph{Level 3: Governance verdict.}
The verdict engine checks:
\begin{itemize}
\item $\OB(\text{pay})$ holds (evidence above threshold).
\item $\PE(\text{pay})$ holds (from \code{ob\_implies\_pe}).
\item $\neg\OB(\neg\text{pay})$ holds (\DTS{} consistency: the agreement
      does not also oblige non-payment).
\item No conflict detected: the \DTS{} trichotomy classifies the tenant
      as obligated, not optional or forbidden.
\end{itemize}
Result: \code{GovernanceVerdict.compliant}.

If the tenant fails to pay, the CTD obligation $\OB(\text{notice} \mid \neg\text{pay})$
activates via the conditional obligation operator (\code{cond\_obl}).
The verdict shifts to \code{violation} for the primary obligation, but the
CTD provides a recovery path --- exactly the scenario \DDL{} was designed to handle.

%% ============================================================
\section{GovNormCycle: Constructive Reactive Accessibility}
\label{sec:govnorm}
%% ============================================================

The \code{ClosedGovAccessibility} structure requires two proofs:
\emph{seriality} (every live state has a successor) and
\emph{closure} (successors of live states are live).
A naive approach using a terminating language such as PyashCore fails:
the liveness predicate ``\code{instr $\neq$ Done}'' is not forward-closed under
reduction, as witnessed by \code{pyashCore\_isGovLive\_not\_closed}.

\paragraph{The GovNormCycle language.}
We provide a purpose-built 2-state reactive loop:

\begin{lstlisting}
def govDeliberateState : Pattern := .apply "GovDeliberate" []
def govEnactState      : Pattern := .apply "GovEnact" []

inductive GovNormCycleStep : Pattern -> Pattern -> Prop where
  | deliberateToEnact : GovNormCycleStep govDeliberateState govEnactState
  | enactToDeliberate : GovNormCycleStep govEnactState govDeliberateState

def govNormLive : Pattern -> Prop :=
  fun p => p = govDeliberateState \/ p = govEnactState
\end{lstlisting}

\noindent
Both states are permanently live.  Seriality and closure follow directly from
the two constructors:

\begin{lstlisting}
theorem govNormCycle_serial :
    forall p, govNormLive p -> exists q, GovNormCycleStep p q := by
  intro p hp
  rcases hp with rfl | rfl
  . exact <|govEnactState, .deliberateToEnact|>
  . exact <|govDeliberateState, .enactToDeliberate|>

theorem govNormCycle_closed :
    forall p q, govNormLive p -> GovNormCycleStep p q -> govNormLive q := by
  intro p q _hp hpq
  cases hpq
  . right; rfl   -- deliberateToEnact: q = govEnactState
  . left; rfl    -- enactToDeliberate: q = govDeliberateState
\end{lstlisting}

\begin{theorem}[Constructive governance accessibility]
The 2-state deliberate-enact cycle satisfies seriality and closure,
providing a constructive witness that \code{ClosedGovAccessibility}
is satisfiable: every live state has a successor, and successors of
live states are live.
\end{theorem}

\noindent
This concrete instance can be composed with any \code{GovFrame} to produce a
\code{GovernanceDDLBundle}.  The model is deliberately minimal: the two states
represent the idealised governance cycle (deliberation alternates with
enactment) without encoding any domain-specific policy content.

\begin{remark}[Why not PyashCore?]
One might expect to instantiate governance accessibility using a richer
operational language such as PyashCore.  However, the liveness predicate
``\code{instr $\neq$ Done}'' is \emph{not} forward-closed under reduction:
PyashCore programs can reach \code{Done}, violating seriality.
This is proven as \code{pyashCore\_isGovLive\_not\_closed}.
The \code{GovNormCycle} is therefore the simplest language that satisfies
the seriality requirement, serving as a proof of satisfiability while
leaving richer instantiations to future work.
\end{remark}

%% ============================================================
\section{Temporal-Deontic Bridge}
\label{sec:temporal}
%% ============================================================

\subsection{Temporal Meanings}

To model time-varying obligations we pair a base world type with a time type:

\begin{lstlisting}
abbrev TemporalWorld (w T : Type*) := w x T
abbrev TemporalMeaning (c w T : Type*) := Meaning c (w x T)
\end{lstlisting}

\noindent
Lead (future) and lag (past) shifts on meanings capture how propositions vary
with time (where \code{T} has \code{[Add T]} or \code{[Sub T]} respectively):

\begin{lstlisting}
-- phi evaluated Delta time units ahead
def meaningLead (phi : TemporalMeaning c w T) (Delta : T) :=
  fun ctx (v, t) => phi ctx (v, t + Delta)

-- phi evaluated Delta time units in the past
def meaningLag (phi : TemporalMeaning c w T) (Delta : T) :=
  fun ctx (v, t) => phi ctx (v, t - Delta)
\end{lstlisting}

\subsection{Temporally Uniform Obligation}

A \code{DDLPlusFrame} on $w \times T$ is \emph{temporally uniform} if the
accessibility and obligation relations are stable under time translation:

\begin{lstlisting}
structure TemporallyUniformOb {w T : Type*} [AddCommGroup T]
    (F : DDLPlusFrame (w x T)) where
  -- pv is translation-invariant
  pv_shift : forall (Delta : T) (v : w) (t : T) (v' : w) (t' : T),
      F.pv (v, t) (v', t') <-> F.pv (v, t + Delta) (v', t' + Delta)
  -- ob is compatible with base-point shift (content lags by Delta)
  ob_shift : forall (Delta : T) (v : w) (t : T) (Y : WProp (w x T)),
      F.ob (F.pv (v, t)) Y <->
        F.ob (F.pv (v, t + Delta)) (fun p => Y (p.1, p.2 - Delta))
\end{lstlisting}

\noindent
Any concrete temporal frame must discharge these shift-invariance conditions.

\subsection{Obligation Persistence Theorem}

\begin{theorem}[Temporal-lag persistence of ideal obligation]
\label{thm:temporal}
Let $F$ be a \code{DDLPlusFrame} on $w \times T$ where $T$ is an abelian group,
and let \code{tub : TemporallyUniformOb F}.  If $\varphi$ is ideally obligatory
at world $(v, t_0)$, then the lag-shifted version $\varphi[\cdot - \Delta]$
is ideally obligatory at $(v, t_0 + \Delta)$.
\end{theorem}

\paragraph{Significance.}
This theorem captures the intuition that obligations persist under
temporal perspective shifts.  For example, if a tenant is obligated to pay
rent on the 1st of the month, then from the perspective of the 15th the
obligation to \emph{have paid} 14 days ago still holds.  More precisely:
obligations do not disappear when one changes temporal reference frame ---
they shift with it, with the content lagging by the same amount as the
reference point advances.

\noindent
\emph{Proof sketch.}
The proof has three parts:
\begin{enumerate}
\item \emph{ob part}: \code{ob\_shift} maps
      $F.\text{ob}(F.\text{pv}\,(v, t_0),\, \varphi)$ to
      $F.\text{ob}(F.\text{pv}\,(v, t_0+\Delta),\, \lambda p.\,\varphi(p_1, p_2{-}\Delta))$,
      which is exactly \code{meaningLag phi Delta}.
\item \emph{pv part}: the violation witness $(v_w, t_w)$ is
      shifted to $(v_w, t_w + \Delta)$ using \code{pv\_shift}.
\item \emph{Negation}: at the shifted witness, \code{add\_sub\_cancel\_right}
      reduces the lag-shift evaluation back to the original negation. \qed
\end{enumerate}

\begin{remark}
The ``lag'' direction is the only direction provable from \code{ob\_shift} as
stated: the frame does not directly relate $\text{ob}$ at two different base
times without the content shift.  Shifting the reference time forward requires
adjusting all content backward to keep the same relative obligations.
\end{remark}

\subsection{Causal-Normative Prohibition}

We connect the temporal-deontic bridge to \PLN{}'s causal heuristics
\cite{goertzel2009pln}.
A \emph{spurious causal candidate} is an element $x$ at time $t$ with
extensional but not intensional predictive support:

\begin{lstlisting}
structure SpuriousCausalCandidate
    (profile : PredictiveCausalProfile Domain Time)
    (x : Domain) (t : Time) : Prop where
  extensional    : profile.extensional x t
  not_intensional : not (profile.intensional x t)
\end{lstlisting}

\begin{theorem}[Spurious prohibition]
If a \DTS{} obligation structure obliges only intensionally-grounded elements
(\code{ObligesOnlyIntensional}), then no spurious causal candidate can be
obligated.
\end{theorem}

\noindent
The proof is a one-line contradiction: \code{hSpur.not\_intensional}
contradicts \code{hGnd.ob\_intensional x hOb t}.
The contrapositive (\code{obligatory\_not\_spurious}) is derived immediately.

%% ============================================================
\section{Abstract Modal Query Encoder}
\label{sec:mqe}
%% ============================================================

The event-calculus \cite{KowalskiSergot1986,Shanahan1999} and the governance
reasoning layer both define query encoders of the form
\code{Modality -> Referent -> Query}, but they were independently defined.
We introduce an abstract structure that unifies both:

\begin{lstlisting}
structure ModalQueryEncoder (Modality Referent Query : Type*) where
  encode : Modality -> Referent -> Query
\end{lstlisting}

\paragraph{Event-calculus instance.}
The event calculus uses three temporal modalities:

\begin{lstlisting}
inductive EventModality where
  | holdsAt      -- fluent holds at time t
  | initiatedAt  -- fluent initiated at t
  | terminatedAt -- fluent terminated at t
\end{lstlisting}

\noindent
An \code{EventQueryEncoder Event Time Query} (with fields \code{holdsAt},
\code{initiatedAt}, \code{terminatedAt}) is equivalent to a
\code{ModalQueryEncoder EventModality (Event x Time) Query}:

\begin{lstlisting}
def eventQueryEncoderToModal (enc : EventQueryEncoder Event Time Query) :
    ModalQueryEncoder EventModality (Event x Time) Query where
  encode
    | .holdsAt,      (e, t) => enc.holdsAt e t
    | .initiatedAt,  (e, t) => enc.initiatedAt e t
    | .terminatedAt, (e, t) => enc.terminatedAt e t
\end{lstlisting}

\begin{theorem}[Event-encoder roundtrip]
The maps \code{eventQueryEncoderToModal} and \code{modalToEventQueryEncoder}
are mutually inverse.
Both roundtrips are proven as simp lemmas using \code{cases modality} and \code{rfl}.
\end{theorem}

\paragraph{Deontic instance.}
The \code{DeonticQueryEncoder.modalQuery} field has exactly the shape
\code{DeonticModality -> Eventuality Entity Pred -> Query}, giving:

\begin{lstlisting}
def deonticQueryEncoderToModal (d : DeonticQueryEncoder Entity Pred Query) :
    ModalQueryEncoder DeonticModality (Eventuality Entity Pred) Query where
  encode := d.modalQuery
\end{lstlisting}

\noindent
The embedding is a direct field projection.
With this unification, dispatch infrastructure, traversal utilities, and
test harnesses written for \code{ModalQueryEncoder} apply to both the
event-calculus and deontic encoders without duplication.

%% ============================================================
\section{Gewirth's Principle of Generic Consistency}
\label{sec:pgc}
%% ============================================================

The Principle of Generic Consistency (\PGC{}) \cite{Gewirth1978} asserts:
\begin{quote}
\emph{Act in accord with the generic rights of your recipients as well as of yourself.}
\end{quote}
Gewirth's argument proceeds from agency: an agent who acts on purpose must
claim the necessary conditions for agency (freedom and well-being, \emph{FWB})
as rights, and by universalizability must grant the same rights to all agents.

Benzm{\"u}ller et al.\ \cite{Benzmuller2019PGC} gave the first machine-checked
formalization of the \PGC{} in Isabelle/HOL, available in the Archive of Formal Proofs.
Our Lean~4 port follows their proof structure while connecting it to the \DDL{}
and \DTS{} layers developed in the preceding sections.

The formalization in \code{Mettapedia/Ethics/} uses the \DDL{} embedding
developed above:

\begin{itemize}
\item \code{Core.lean}: introduces the agent world type, purpose-directed
      action predicate \code{PPA}, and the generic rights type \code{FWBRight}.
\item \code{GewirthPGC.lean}: packages a \code{PGCInterpretation} (frame +
      assignments) and states the main theorem:
      \code{PGC\_strong\_ofAssumptions} proves that any agent satisfying
      \code{PGCAssumptions} is ideally obligated to respect the FWB rights
      of all purposive agents.
\item \code{GewirthBridge.lean}: connects the local \code{Oi} and \code{Ocond}
      operators in \code{GewirthPGC.lean} to the corresponding \code{ideal\_obl}
      and \code{cond\_obl} from \code{DDLPlus.Core}, making the bridge explicit.
\end{itemize}

\noindent
The \DTS{} interdefinability theorems (§\ref{sec:dts}) underlie the Gewirth
formalization: the equivalences between obligation, permission, and optionality
ensure that the logical structure of the PGC proof is sound regardless of which
representation is chosen.

%% ============================================================
\section{MeTTa-to-Lean Refinement Theorems}
\label{sec:refinement}
%% ============================================================

The \code{governance-reasoning-engine} \cite{GovReasoningEngine} implements
\DTS{} rules as MeTTa rewrite rules.  A natural question is whether these
operational rules \emph{correctly encode} the \DTS{} theorems proven
semantically in Lean.  This section presents formal refinement theorems
answering this question for three evaluator levels: PeTTa (type-free),
HE~MeTTa (with type and binding threading), and an answer-level projection
that abstracts over both.

\subsection{Architecture}

The refinement bridge has four layers:

\begin{enumerate}
\item \textbf{Rule encoding.} Each MeTTa governance rule (e.g.,
      \code{(= (ct-triple-for-add \$e type permitted) (ct-triple \$e type obligatory))})
      is encoded as a concrete Lean \code{RewriteRule} value with
      \code{premises = []}.
\item \textbf{Deontic interpretation.} A function
      \code{interpretDeontic : Pattern -> DeonticInterp}
      classifies ground patterns into deontic assertions:
      \code{ct-triple} patterns map to \code{.modalAssertion e .obligatory},
      \code{ct-triple-for-add} patterns to \code{.modalAssertion e .permitted},
      and so on for \code{.forbidden} and \code{.optional}.
\item \textbf{Interpretation correctness.} For each rule, we prove that
      the left-hand side and right-hand side, after binding substitution,
      interpret to the expected deontic modalities.
\item \textbf{Semantic refinement.} We bundle the interpretation correctness
      proofs with the corresponding \DTS{} theorem, showing that the rule
      structure matches the proven algebraic property.
\end{enumerate}

\subsection{PeTTa Refinement}

PeTTa is the type-free, binding-free evaluator: \code{PeTTaEval s p [q]} states
that pattern \code{p} evaluates to \code{[q]} in space \code{s} via rewrite
rule application.

\begin{theorem}[PeTTa OB$\Rightarrow$PE refinement]
\label{thm:petta-ob-pe}
For any string \code{e}:
\begin{enumerate}
\item $\mathit{interpretDeontic}(\mathit{applyBindings}\;[(\text{``e''}, e)]\;\mathit{rule.right})
       = \mathit{modalAssertion}\; e\; \mathit{obligatory}$,
\item $\mathit{interpretDeontic}(\mathit{applyBindings}\;[(\text{``e''}, e)]\;\mathit{rule.left})
       = \mathit{modalAssertion}\; e\; \mathit{permitted}$, and
\item $\forall\, (d : \DTS{}\;P)\;(p : P),\; d.\mathit{ob}\;p \to d.\mathit{pe}\;p$.
\end{enumerate}
\end{theorem}

\noindent
Parts~1 and~2 are proved by \code{simp} on the concrete rule encoding and
binding substitution.  Part~3 is the already-proven \DTS{} theorem
\code{ob\_implies\_pe}.  The three-part conjunction bundles rule-level
pattern interpretation with the semantic theorem it encodes.

Analogous refinement theorems are proven for $\OB(\neg p) \Rightarrow \neg\PE(p)$
(\code{dts\_ob\_neg\_not\_pe\_refinement}), $\neg\PE(p) \Rightarrow \OB(\neg p)$,
$\OP(p) \Leftrightarrow \neg\OB(p) \wedge \neg\OB(\neg p)$,
$\OB(p) \Rightarrow \neg\OP(p)$, and the \DTS{} trichotomy --- all 7 core
interdefinability relations from Table~\ref{tab:dts}.

\subsection{Shared MeTTaLike Interface and Let* Equivalence}

The MeTTa \code{let*} construct is not a built-in primitive but a pair of
rewrite rules (\code{letStarBaseRule}, \code{letStarRecRule}) that unfold
\code{let*} to nested \code{let} expressions.  To prove these unfoldings
generically, we introduce a typeclass:

\begin{lstlisting}
class MeTTaLike (Eval : PeTTaSpace -> Pattern -> List Pattern -> Prop) where
  ruleApp : forall {s r bs p},
    r in s.rules -> r.premises = [] -> bs in matchPattern r.left p ->
    Eval s p [applyBindings bs r.right]
\end{lstlisting}

\noindent
Both \code{PeTTaEval} (the type-free evaluator) and \code{HEEvalAnswers}
(the answer-level projection of HE~MeTTa) are instances.
With this abstraction, unfolding theorems hold for \emph{any} MeTTaLike evaluator:

\begin{theorem}[Let* unfolding]
For any MeTTaLike evaluator \code{Eval} and space \code{s} containing the
\code{let*} rules:
\begin{align*}
  &\mathit{Eval}\;s\;(\mathit{let*}\;[(v_1, e_1), (v_2, e_2)]\;\mathit{body})
    \;[\mathit{let}\;v_1\;e_1\;(\mathit{let*}\;[(v_2, e_2)]\;\mathit{body})] \\
  &\mathit{Eval}\;s\;(\mathit{let*}\;[(v_2, e_2)]\;\mathit{body})
    \;[\mathit{let}\;v_2\;e_2\;(\mathit{let*}\;[]\;\mathit{body})] \\
  &\mathit{Eval}\;s\;(\mathit{let*}\;[]\;\mathit{body})\;[\mathit{body}]
\end{align*}
\end{theorem}

\noindent
The match proofs use \code{simp} on the concrete \code{matchBag} semantics
for bag patterns with rest variables.  The full unfolding sequence for
2-binding and 3-binding \code{let*} chains is bundled as
\code{letStar\_full\_2} and \code{letStar\_full\_3}.
This is significant because the DTS governance rules in
\code{DTS.metta} (lines 14--135) use \code{let*} chains extensively;
the unfolding theorems ensure these chains reduce correctly under
any compliant evaluator.

\subsection{HE MeTTa Refinement}

The HE (Hyperon Experimental) MeTTa evaluator extends PeTTa with explicit
type annotations and binding threading:
\code{MeTTaEval s p ty inputBindings results} produces result pairs
\code{(pattern, bindings)}.  The answer-level projection erases
types and bindings:

\begin{lstlisting}
def HEEvalAnswers (s : PeTTaSpace) (p : Pattern)
    (answers : List Pattern) : Prop :=
  exists ty inputBs results,
    MeTTaEval s p ty inputBs results
    /\ results.map Prod.fst = answers
\end{lstlisting}

\noindent
All PeTTa refinement theorems lift automatically to \code{HEEvalAnswers}
via the shared \code{MeTTaLike} instance.  Additionally, full
\code{MeTTaEval}-level theorems are proven with explicit binding tracking:

\begin{theorem}[HE binding threading for OB$\Rightarrow$PE]
If \code{dtsRule\_ob\_implies\_pe} is in a PeTTa space, the HE~MeTTa evaluator
fires it with match bindings \code{[("e", .apply e [])]} prepended to the
input bindings, and the interpretation/semantic refinement hold as before.
\end{theorem}

\subsection{Canonical HE Spec Architecture}

The Lean formalization contains two representations of HE~MeTTa semantics:

\begin{itemize}
\item \textbf{Computable interpreter}: \code{Languages/MeTTa/HE/Interpreter.lean}
      defines 6 mutually recursive functions (\code{metta},
      \code{interpretExpression}, \code{interpretFunction},
      \code{interpretArgs}, \code{interpretTuple}, \code{mettaCall})
      with fuel-bounded execution, matching the Rust reference implementation.
      A suite of 37 conformance theorems
      (\code{HE/Conformance.lean}) verifies this interpreter against the spec.
\item \textbf{Inductive relation}: \code{OSLF/PeTTa/MeTTaEval.lean} provides
      the same semantics as an inductive proposition, amenable to
      proof-theoretic reasoning.  The refinement theorems in this section
      use this representation.
\end{itemize}

\noindent
Bridging the computable interpreter to the inductive relation requires
reconciling different type representations (\code{Atom} with 4~constructors
vs.\ \code{Pattern} with 7~constructors) and different binding models
(\code{HE.Bindings} with assignments and equalities vs.\ simple association
lists).  A lossy bridge (\code{patternToAtom}/\code{atomToPattern}) exists
in \code{MeTTaCore/Bridge.lean}; a full general equivalence proof remains
future work, but the DTS rule subset is \emph{exact}: governance rules
use only \code{apply} and \code{fvar} pattern constructors, which map
cleanly to \code{expression} and \code{var} atoms.

For this exact subset, \code{HECanonicalConformance.lean} proves that the
canonical HE interpreter produces the correct results via kernel-checked
\code{rfl} proofs:

\begin{lstlisting}
theorem he_canonical_ob_pe_soaMoor :
    eval (expression [symbol "ct-triple-for-add",
            expression [symbol "soaMoor"], ...])
         dtsHESpace =
    [(expression [symbol "ct-triple",
            expression [symbol "soaMoor"], ...],
      { assignments := [("e", ...)], equalities := [] })] := rfl
\end{lstlisting}

\noindent
The Lean kernel itself performs the full HE evaluation (through \code{metta}
$\to$ \code{interpretExpression} $\to$ \code{mettaCall} $\to$
\code{queryEquations} $\to$ \code{simpleMatch} $\to$ \code{Bindings.apply})
and verifies the output.  End-to-end conformance bundles eval, decode
(\code{atomToPattern}), deontic interpretation, and the DTS theorem into
a single conjunction --- all kernel-checked.

\subsection{What the Refinement Proves}

The refinement theorems establish:

\begin{quote}
\emph{If the governance-reasoning-engine DTS rules are encoded as PeTTa rewrite
rules, then evaluating them produces results that, under deontic interpretation,
satisfy the corresponding \DTS{} theorems proven in Lean.}
\end{quote}

\noindent
This is a soundness result at three levels:

\begin{enumerate}
\item \textbf{Inductive semantics}: The \code{PeTTaEval} and
      \code{MeTTaEval} relations fire the DTS rules with correct
      binding substitution and deontic interpretation (7 rules each).
\item \textbf{Shared evaluator abstraction}: The \code{MeTTaLike}
      typeclass ensures \code{let*} unfolding works identically for
      both evaluators (2- and 3-binding chains proven).
\item \textbf{Canonical interpreter}: The computable HE \code{eval}
      function produces the correct results for all three DTS rules,
      verified by \code{rfl} (the Lean kernel performs the computation).
\end{enumerate}

\noindent
The verification does not cover the full MeTTa codebase, but it does
verify that the core deontic rewrite rules faithfully encode the
algebraic properties they claim to implement, and that the canonical
interpreter agrees with the inductive semantics on the DTS rule subset.

All refinement proofs are kernel-checked with zero \code{sorry} placeholders
and use only the standard axioms (\code{propext}, \code{Classical.choice},
\code{Quot.sound}).

%% ============================================================
\section{Related Work}
\label{sec:related}
%% ============================================================

\paragraph{Isabelle/HOL formalizations.}
Benzm{\"u}ller et al.\ \cite{Benzmuller2018} formalized \DDL{} in Isabelle/HOL
using the same shallow semantic embedding approach; our Lean~4 \DDL{} formalization
is a port of that work.
Benzm{\"u}ller and Fuenmayor \cite{Benzmuller2019PGC} formalized the Gewirth \PGC{}
in Isabelle/HOL, available in the Archive of Formal Proofs; our \PGC{}
formalization is likewise a port.
We extend both with the governance reasoning layer, temporal bridges,
and \PLN{} causal integration.
The Isabelle sources (\code{CJDDLplus.thy}, \code{ExtendedDDL.thy}) are
included in the repository as a reference.

\paragraph{Input/Output logic.}
Makinson and van der Torre's input/output logic \cite{MakinstonVanDerTorre2000}
provides an alternative semantics for conditional obligation.  The \code{ob}
dyadic relation in \DDL{} can be interpreted in their framework, though the
present formalization does not develop this connection explicitly.

\paragraph{MeTTa governance reasoning engine.}
The \code{governance-reasoning-engine} \cite{GovReasoningEngine} implements
a three-level inference pipeline (eventuality, statement, verdict) in MeTTa,
using non-deterministic evaluation for governance verdict computation.
Our formalization verifies the algebraic properties (\DTS{}
interdefinability, evidence monotonicity) that correspond to structures
used in their implementation, and the refinement theorems of
\S\ref{sec:refinement} formally connect the MeTTa rule encodings to the
Lean semantics for all 7 core \DTS{} interdefinability relations.

\paragraph{Event calculus in Lean.}
The \PLN{} Probabilistic Event Calculus formalized in
\code{PLNProbabilisticEventCalculus.lean} (this repository) follows the
Kowalski-Sergot \cite{KowalskiSergot1986} and Shanahan \cite{Shanahan1999}
presentations.  The \code{ModalQueryEncoder} of §\ref{sec:mqe} unifies this
with the deontic layer.

%% ============================================================
\section{Limitations}
\label{sec:limitations}
%% ============================================================

\begin{itemize}
\item \textbf{No end-to-end regulatory corpus.}
      The formalization provides verified logical primitives but has not been
      tested against any regulatory corpus (GDPR, SEC rules, etc.).
\item \textbf{No natural language parser.}
      There is no automated pipeline from regulatory text to \DDL{} frame
      instances; all examples are hand-encoded.
\item \textbf{No benchmark suite.}
      We do not provide a test suite of governance scenarios with expected
      verdicts against which implementations could be evaluated.
\item \textbf{Rule-level (not full-program) refinement.}
      The refinement theorems (\S\ref{sec:refinement}) verify that individual
      DTS rewrite rules correctly encode the corresponding \DTS{} theorems,
      but they do not verify the full MeTTa governance-reasoning-engine
      codebase.  The three-level inference pipeline (eventuality, statement,
      verdict) is verified at the rule-firing level, not end-to-end.
      The canonical HE interpreter and the inductive \code{MeTTaEval}
      relation agree on the DTS rule subset (proven via \code{rfl}
      conformance tests), but a general equivalence has not been established.
\item \textbf{No deployed case study.}
      The system has not been deployed in any real governance or compliance
      setting.
\item \textbf{DTS consistency requires external hypothesis.}
      The \DTS{} consistency axiom ($\OB(p) \to \neg\OB(\neg p)$) is not
      derivable from \DDL{} frame conditions alone; the bridge
      \code{frameToDTS} requires an explicit \code{hcons} parameter
      (see Remark in §\ref{sec:dts}).
\item \textbf{Minimal accessibility model.}
      The \code{GovNormCycle} is a 2-state proof of satisfiability, not a
      practical governance state machine.
\end{itemize}

%% ============================================================
\section{Future Developments: Toward Practical Governance}
\label{sec:future}
%% ============================================================

The formalization presented here establishes a verified logical foundation
but does not yet constitute a deployable governance tool.  Several
real-world systems demonstrate that deontic logic \emph{can} deliver
practical value, and point toward concrete development paths.

\paragraph{Regulatory compliance checking.}
Governatori and Shek's \emph{Regorous} system \cite{Governatori2013}
uses defeasible deontic logic to check business processes (modelled in
BPMN) against regulatory requirements, including GDPR formalization.
A natural extension of our work would be to compile \DDL{} frame
instances from regulatory texts (using NLP extraction or GF grammars)
and check process traces against the resulting obligation/permission
structure.  The \DTS{} trichotomy (obligatory/optional/forbidden) maps
directly to compliance verdicts.

\paragraph{Multi-agent governance for AI systems.}
XMPro's Multi-Agent Generative Systems (MAGS) platform \cite{XMProMAGS}
uses deontic principles as ``Rules of Engagement'' for industrial AI agents:
obligations that agents must fulfil, permissions that bound their autonomy,
and prohibitions they cannot violate.  These are encoded as \emph{deontic tokens}
embedded in agent profiles.  Our \DDL{} formalization could provide verified
backing for such tokens: each deontic token would carry a proof witness
(from the CJ theorems) ensuring it is consistent with the frame conditions,
ruling out contradictory obligations before deployment.

\paragraph{AI-to-AI treaty layer.}
The minimal \code{TreatyKernel} introduced here points toward a more general
treaty layer for decentralized agent economies: machine-checkable bilateral or
multilateral agreements specifying delegated tasks, deadlines, privacy
constraints, payment conditions, and contrary-to-duty repair chains.  In such a
setting, abstract norms should be checked against concrete runtime traces via
subsumption rather than exact event identity.  The next step is to extend the
kernel with activation conditions, repair discharge, delegation scope, and
compositional reasoning over networks of treaties.

\paragraph{Contrary-to-duty policy repair.}
The Chisholm paradox (§\ref{sec:rental}) shows that naive deontic logic
produces contradictions when primary obligations are violated.  \DDL{}'s
dyadic obligation resolves this, but the formalization does not yet support
automated \emph{repair}: given a detected violation, synthesize the
applicable contrary-to-duty obligations and present them as a recovery
plan.  This would require integrating the temporal bridge (§\ref{sec:temporal})
with a planning layer that enumerates CTD chains.

\paragraph{Richer governance accessibility models.}
The \code{GovNormCycle} (§\ref{sec:govnorm}) is a minimal 2-state
witness of satisfiability.  A practical system needs governance
accessibility over richer state spaces --- for instance, a reactive
event-loop language where each ``state'' is a snapshot of active
norms, pending obligations, and violation history.  The negative
result for PyashCore (\code{pyashCore\_isGovLive\_not\_closed}) shows
that terminating languages fail the seriality requirement; future
work should explore co-inductive or stream-based state machines
that guarantee perpetual liveness.

\paragraph{Probabilistic evidence integration.}
The three-level judgment pipeline (§\ref{sec:govlayer}) uses
\PLN{} evidence values but does not yet connect them to posterior-mean
inference.  Integrating the \DTS{} strength values with the
probabilistic event calculus (via the \code{ModalQueryEncoder}) would
allow governance verdicts to carry calibrated confidence levels,
enabling graded compliance rather than binary pass/fail.

\paragraph{Occurrence semantics beyond \code{rexist}.}
A more temporally explicit alternative to the Hobbs-style \code{rexist}
predicate is to model actuality through event occurrence and fluent
persistence, in the style of the Event Calculus~\cite{KowalskiSergot1986}:
occurrence at time~$t$, together with initiates/terminates relations, yields
time-indexed holding conditions rather than a timeless actuality marker.  The
present treaty kernel already provides a lightweight step in this direction
via timestamped treaty events and deadline reasoning (the \code{occursAt}
predicate in \code{OccurrenceMVP.lean}).  A natural extension is to replace
timeless \code{rexist}-queries with policy-filtered accepted occurrences and
then lift treaty reasoning to time-indexed fluents such as escrow status,
data-retention status, and breach status.

\paragraph{Explicit reasons and provenance.}
Another extension is to integrate justification-style reason
terms into the governance/world-model interface.  Instead of
representing only that an eventuality is obligatory, forbidden, or
accepted as having occurred, one could represent the explicit
\emph{reasons} supporting these statuses: attestations, sensor reports,
derivation rules, or policy certificates.  The \code{attestedBy} field
on \code{TreatyEvent} already exposes a provenance hook; a dedicated
\emph{attestation query sort} would lift provenance into the typed
world-model interface.  This connects naturally to W3C PROV and to
justification logic, where modal claims are unfolded into explicit
reason terms $t : \varphi$.

\paragraph{Typed governance query sorts.}
The current \code{DeonticQueryEncoder} distinguishes two query channels
(modal/eventuality and ground/triple).  A natural refinement is to
introduce explicit governance query sorts --- \emph{norm}, \emph{occurrence},
\emph{projection}, \emph{attestation} --- as a sort-indexed typed layer
over the abstract world-model interface.  The existing \code{WorldModelSigma}
typed query infrastructure already supports this without modifying the
untyped WM calculus core.

%% ============================================================
\section{Conclusion}
\label{sec:conclusion}
%% ============================================================

We have presented a Lean~4 formalization of governance deontic logic
spanning six layers: \DDL{} semantic embedding, \DTS{} algebra, Hobbs
eventuality reification with governance judgments, constructive reactive
accessibility, temporal-deontic and causal-normative bridges, the
Gewirth \PGC{}, and formal MeTTa-to-Lean refinement theorems.

The refinement chain is the deepest contribution: it connects
operational MeTTa governance rules to proven \DTS{} theorems at three
levels --- inductive semantics (PeTTa and HE~MeTTa), a shared evaluator
abstraction with \code{let*} unfolding, and kernel-checked canonical
interpreter conformance.  For the DTS rule subset, the Lean kernel
itself performs the HE evaluation and verifies the output.

The addition of event subsumption and a minimal treaty kernel moves the work
from static governance verdicts toward proof-carrying treaty execution over
concrete agent interaction traces.  The evidence-graded actuality layer
(§\ref{sec:actuality}) replaces the metaphysical ``really exists'' reading
of \code{rexist} with a two-layer semantics --- graded eventuality
evidence and policy-relative categorical acceptance --- and the
provenance-aware accepted-trace mechanism demonstrates the full pipeline
from evidence assessment through attestor filtering to treaty compliance.

The formalization is compositional (each layer is independently usable)
and grounded in published sources (Carmo \& Jones 2002, von Wright 1951,
Hobbs 1985, ISO~24617-4).

%% ============================================================
\bibliographystyle{plain}
\bibliography{references}
%% ============================================================

\end{document}
